\subsection{Integral transfer from the tensor of two-letter spaces}
\label{sec:edge-transfer}
\hypertarget{source-edge-transfer}{}

Each top in the tensor model has independent vector and covector variables. Transferring its tests to the actual weight space requires an integral comparison with the irreducible summands.

Let \(T\) be the set of tops on the cycle and \(O\) the coefficient valuation ring. Put \(N=|T|\), \(k=m+1\), and \(K=\operatorname{Frac}(O)\).
Use \(A_s,B_s\) for the two nonnegative degrees at a top.
The tensor model \(U\) takes independent vector and covector
variables at each top. At an odd top their roles switch, so the
first-coordinate normalization of \(U\) is
\[
 a_U=\sum_{s\ {\rm even}}A_s+\sum_{s\ {\rm odd}}B_s.
\]
Use the component frames of the preceding subsection for the
coefficient lattice \(U(M)\), for a tuple \(M\) of \(O\)-lattices.
Apply the definitions of the edge tests to \(U\) with \(a_U\) in
place of \(A_1\), including the pairing \(p^{a_U}\phi(u)\).

In one component let \(P_{A,B}\) be the polynomials of bidegree
\((A,B)\) in \(x,y\), and let \(h=\sum_i x_i y_i\) and
\[
 \Delta=\sum_i\partial_{x_i}\partial_{y_i}.
\]

Writing \(\beta_{A,B}\) for the quotient by \(hP_{A-1,B-1}\),
this quotient is the dual Weyl model for
\((A,0,\ldots,0,-B)\), with its determinant twist.
The identification holds on coefficient lattices. To see the
normalization, send \(x\) to the first column and \(y\) to the
signed \((k-1)\)-row minors of the first \(k-1\) columns. Use
independent matrices for the two-letter space and one matrix
for the quotient. Under a change of the row frame, the column
transforms as a vector and the minors as a covector times the
determinant. The twist cancels the determinant.

\begin{proposition}[Transfer of a tensor test to the weight summands]\label{prop:edge-test-transfer}
Take \(k\ge3\), \(w^2=p\), and \(A_s+B_s+1\le8a\) at each top.
Put \(u=N\lceil(k+1)\log_p(8a+k-1)\rceil\). Suppose a test for
\(U\) has \(r_U\) directions at depth \(D_0\) at a fixed edge,
with the specified endpoint order and conductor. Then there are
tests with \(r_j\) directions on the summands \(V_j\) of weights
\((A_s-j_s,0,\ldots,0,-B_s+j_s)\),
\(0\le j_s\le\min(A_s,B_s)\), at depth \(D_0+3u\), with
the same endpoint order and conductor and total rank bounded below by
\[
 \sum_j r_j\ge r_U,\qquad D_0+3u\le D_1=D_0+4u.
\]
In particular the tests are available at depth \(D_1\).
\end{proposition}

\begin{proof}
\proofstep{The injective substitution and differential operators.}
In the independent model the substitution is injective over
every field, since each nonzero covector is the exterior
product of some frame of its kernel, with one frame variable
rescaled to give the specified exterior product. In the
common-matrix model the image is dense in the incidence
hypersurface \(h=0\). Over an algebraic closure, for any
nonzero \(x,y\) with \(y(x)=0\), extend \(x\) to a basis of
\(\ker(y)\) and adjust one of the other vectors to realize
\(y\). Here \(k\ge3\). The polynomial \(h\) is irreducible, by
viewing it as primitive of degree one in any \(y_i\); hence the ideal of relations in the common-matrix substitution is
\[
 \ker(\text{common-matrix substitution})=(h).
\]
These arguments also apply over the residue
field after base change to its algebraic closure.

Over \(O\)
one may divide by the monic leading monomial of \(h\), so the
bihomogeneous quotient is free. The injection after reduction
modulo \(w\) proves saturation of its image. In the standard
frame the integral polynomial realization of the dual Weyl
module is its span over \(K\) intersected with the coefficient
lattice. Equivalently one uses the integral standard monomial
basis in products of column minors \cite[Theorem II.2.16]{abw-schur}.
The Cartan product of the
first-column and cofactor weights generates that module, giving
the quotient identification. A change of frame gives the assertion for every \(M\).

\proofstep{Projectors on fixed bidegree.}
For \(a,b\ge0\), let \(H_{a,b}\) be the projector onto
\(\ker\Delta\) along \(hP_{a-1,b-1}\).
The commutator formula
\[
 \Delta(h^i u)=h^i\Delta u+i(k+c+d+i-1)h^{i-1}u
\]

for \(u\in P_{c,d}\) gives
\[
 H_{a,b}=\sum_{i=0}^{\min(a,b)}
    \frac{(-1)^i h^i\Delta^i}{i!(k+a+b-2)^{\underline i}},
 \qquad t^{\underline i}=t(t-1)\cdots(t-i+1).
\]
Indeed \(\Delta H=0\), and the constant summand gives the
identity modulo \(h\). If \(u=h^j v\) with \(\Delta v=0\),
\(\Delta^j u=j!(k+A+B-2j)^{\overline j}v\), where the bar
denotes a rising factorial. Induction on \(\min(A,B)\) gives
the unique decomposition into the \(h^j\ker\Delta\).
For \(a=A-j,b=B-j\), its projection and inclusion in
quotient coordinates are
\[
 \pi_j f=\beta_{a,b}\frac{H_{a,b}\Delta^j f}
                        {j!(k+a+b)^{\overline j}},
 \qquad
 \iota_j\beta_{a,b}(u)=h^jH_{a,b}u.
\]
The formula for \(\pi_j\) follows by applying \(\Delta^j\)
to each harmonic summand and killing the terms still divisible
by \(h\).

Both maps have denominator at most
\(p^{(k+1)\lfloor\log_p(A+B+k)\rfloor}\) in every component
frame. For \(\pi_j\), in the \(i\)-th summand of its formula
expand \(\Delta^{i+j}\) by multi-indices \(\nu\) with
\(|\nu|=i+j\). Apart from integral binomial coefficients
from the two differentiated monomials, the coefficient is,
up to sign,
\[
 \frac{\binom{i+j}{i}^2}
 {\binom{i+j}{\nu}\binom{k+a+b-2}{i}\binom{k+a+b+j-1}{j}}.
\]
The factorial formula for a valuation sums one term per \(p\)-power.
For the multinomial that term is at most \(k-1\); each binomial
contributes at most one. All factorial arguments are at most
\(A+B+k\). For \(\iota_j\) the projector formula has the same
estimate with \(j=0\), and an integral lift of the quotient
exists by its coefficient-lattice identification. Both maps
commute with the paired vector and covector action, so one
bound works for all frames.

\proofstep{Tensoring with control of denominators.}
Tensor these maps over the tops. Let \(j=(j_s)_s\),
\(0\le j_s\le\min(A_s,B_s)\), and let \(V_j\) be the weight
space with endpoints \(A_s-j_s,-(B_s-j_s)\). Write
\[
 \pi_j:U\to V_j,\ \iota_j:V_j\to U
\]
for the products of
the maps just constructed. Factoring the sum independently at each top gives
\[
 \sum_{\mathbf j}\iota_{\mathbf j}\pi_{\mathbf j}
 =\bigotimes_s\left(\sum_{j_s=0}^{\min(A_s,B_s)}
        \iota_{s,j_s}\pi_{s,j_s}\right)
 =\bigotimes_s 1=1_U ,
\]
where the map with first index \(s\) is the corresponding single-top
map.

When \(A_s+B_s+1\le8a\), the common lattice allowance is
\[
 u=N\left\lceil (k+1)\log_p(8a+k-1)\right\rceil.
\]
Thus \(p^u\pi_j U(M)\subset H_j(M)\) and
\(p^u\iota_j H_j(M)\subset U(M)\). Put \(a_j=A_1(V_j)\)
and \(\delta_j=a_U-a_j=\sum_s j_s\).
The maps with the normalizations of the edge tests are
\[
 \widetilde\pi_j=w^{\delta_j}\pi_j,\qquad
 \widetilde\iota_j=w^{-\delta_j}\iota_j .
\]
If \(x,\phi\) satisfy the two tests at an edge for \(U\), then
\begin{equation}\label{eq:edge-test-maps}
 p^u\widetilde\pi_j x\in C_e(V_j),\qquad
 p^{2u+\delta_j}\phi\circ\widetilde\iota_j\in C'_e(V_j).
\end{equation}
For source vectors the factor \(w^{\delta_j}\) changes the
edge scale \(w^{-a_U}\) into \(w^{-a_j}\), including at every
other outgoing edge.

At the source itself it only increases
integrality. For covectors use
\(p^u(\phi_0\circ\iota_j)\in H_j(M)^\vee\) whenever
\(\phi_0\in U(M)^\vee\), which follows by evaluating against
\(p^u\iota_j H_j(M)\). The uniform covector estimate in this argument allows
\(2u\) in this step. The factors
\(p^{\delta_j}w^{-\delta_j}=w^{\delta_j}\) give the same
edge and endpoint scales on the dual side. In the normalized
pairing the total effect is exactly \(p^{3u}\), since
\[
 p^{a_j}p^{2u+\delta_j}p^u\,
 \phi(\widetilde\iota_j\widetilde\pi_j x)
       =p^{3u}\,p^{a_U}\phi(\iota_j\pi_j x).
\]
Thus the strength of the extra vanishing at the other edges
is preserved. The exponent \(\delta_j\) may grow with the
weights and is canceled in the pairing.

Let \(zI_{r_U}\) be the matrix of the tensor test after its
integral right multiplication. Apply \eqref{eq:edge-test-maps}
and use the same right multiplier on every summand. This produces
integral matrices \(M_j\) of size \(r_U\) with
\[
                    \sum_j M_j=z_*I_{r_U},
             \qquad z_*=p^{3u}z\ne0 .
\]
For each \(j\), let \(r_j\) count the Smith diagonal entries
dividing \(z_*\), and let \(E_j\subset O^{r_U}\) be the span
of the corresponding left Smith vectors. Over the DVR all
other entries are multiples of \(wz_*\). Consequently the images
of the \(M_j\) in \(O^{r_U}/\sum_j E_j\) are contained in
\(w\) times the image of \(z_*O^{r_U}\). Their sum contains
the image of \(z_*O^{r_U}\), so that image is zero by Nakayama's
lemma. This gives
\[
  r_U=\operatorname{rk}_O\bigl(\textstyle\sum_j E_j\bigr)
                                       \le\textstyle\sum_j r_j .
\]

In the \(j\)-th pairing, take the retained Smith rows and columns
and multiply the right partners by \(z_*\) divided by their
diagonal entries. These factors lie in \(O\), so they preserve
all source and target membership and vanishing tests. The
resulting scalar identity has depth at most \(D_0+3u\).

\end{proof}

\subsection{Combining segment tests at one edge}

Transfer identifies the tensor model on one edge. The segment ranks multiply there only after their pairings and the associated denominator bounds have been combined.

\begin{proposition}[Assembly of the segment tests]\label{prop:segment-edge-assembly}
Use the two-letter tensor \(U\), the retained segment lattices
of ranks \(R_I\), and the segment tests of depths \(D_I\) and
ranks \(r_I\ge(1-\kappa_I)R_I\). Assume the supported
block-pairing congruence and order \(\lfloor t_0(A,B)\rfloor\)
from Sections~\ref{sec:block-pairing}--\ref{sec:symbol-correction},
and let \(D_0=\sum_I D_I+4\kappa_0\). If \(c>2D_0\), the
tensor edge test has depth \(D_0\) and at least \(\prod_I r_I\)
directions. Writing \(r_{\rm edge}\) for the rank of this test, the retained-product bound \eqref{eq:retained-product} gives
\[
 r_{\rm edge}\ge\prod_I r_I\ge\left(1-\kappa_P-\sum_I\kappa_I\right)\frac{\dim U}{g}.
\]

\end{proposition}

\begin{proof}
Choose Smith bases separately on each retained digit block. A
fractional pairing may first be multiplied by a uniformizer
power, which is then removed from the resulting diagonal entries.
If the Smith entries are \(s_i\), the low-valuation rank is
\[
 r_I=\#\{i:v(s_i)\le D_I\},\qquad v(0)=+\infty.
\]

The \(r_I\) left basis vectors span a saturated sublattice; a
primitive vector in it has a unit coordinate, so pairing against
the corresponding right basis vector gives a value of valuation
at most \(D_I\). Conversely, a saturated sublattice of larger
rank has a nonzero intersection with the rational span of the
remaining left vectors. Rescale a vector in the intersection to
be primitive in that sublattice. Every integral partner gives
valuation greater than \(D_I\). Thus the Smith count is exactly
the rank tested by primitive vectors.

Multiply the right partner for a retained entry \(s_i\) by
\[
 w^{\lfloor2D_I\rfloor}/s_i.
\]
The ratio is integral because
\(v(w)=1/2\) and \(2v(s_i)\) is an integer. We obtain diagonal
segment matrices with common value \(w^{\lfloor2D_I\rfloor}\)
and retain the original block support.

\proofstep{The rank after tensoring the segment tests.}
Let \(R_I\) be the rank of the retained lattice for segment \(I\)
and
\[
 r_I\ge(1-\kappa_I)R_I
\]
its successful rank at depth
\(D_I\). Sections~\ref{sec:block-pairing}--\ref{sec:symbol-correction} give the following identity
for a tuple of supported vectors and a tuple of supported
partners. If the support labels are \(\boldsymbol\beta\) and
\(\boldsymbol\beta'\), and the segment pairings are \(a_I\),
the normalized edge pairing is congruent to
\[
  {\bf 1}_{\boldsymbol\beta=\boldsymbol\beta'}\,
                z_{\boldsymbol\beta}\prod_I a_I \pmod{w^c},
   \qquad z_{\boldsymbol\beta}\ne0,\quad
                      v(z_{\boldsymbol\beta})\le4\kappa_0 .
\]
The lifted vectors must satisfy the tests \(C_e(U),C'_e(U)\)
at order \(\lfloor t_0(A,B)\rfloor\) and conductor \(c\).
 The scalar is allowed
to depend on the tuple of blocks. Because the segment
matrices are diagonal, the lifted matrix is diagonal modulo
\(w^c\), with nonzero diagonal values of valuation at most
\(D_0=\sum_I D_I+4\kappa_0\).

If \(c>2D_0\), write the matrix as the diagonal matrix times
an integral matrix congruent to the identity modulo \(w\).
The inverse of the latter is integral. Scaling the
diagonal values to \(w^{\lfloor 2D_0\rfloor}\) then gives a
tensor edge test on \(\prod_I r_I\) directions.

Equation~\eqref{eq:retained-product} supplies the ordinary rank count
on retained blocks.
Combining this with the tests gives
\[
 r_U\ge \left(1-\kappa_P-\sum_I\kappa_I\right)\dim U/g .
\]
In the nonvacuous case all the discarded fractions and their
sum are in \([0,1]\), so the inequality follows by multiplying
the rank bounds and using \(\prod(1-\kappa_I)\ge1-\sum\kappa_I\).
When the right side is nonpositive the zero test suffices.
This proves the required tensor rank at the chosen filtration.

\end{proof}

\subsection{Averaging the rank deficits}

The tensor space decomposes into summands of different dimensions. An estimate for the proportion discarded by the edge test must therefore average their deficits with dimension weights.

Let \(D(a)=\binom{a+k-1}{k-1}\) for \(a\ge0\), set
\(D(-1)=0\), and put
\(d(a,b)=D(a)D(b)-D(a-1)D(b-1)\).
The dimensions are \[
 \dim U=\prod_s D(A_s)D(B_s),\qquad \dim V_j=\prod_s d(A_s-j_s,B_s-j_s).
\]
Summing the latter over \(j\) telescopes to \(\dim U\).
The maximal successful rank at a fixed depth and order is
the same at every edge because \(\mathrm{GL}_k(\operatorname{Frac}W)\)
is transitive on these pairs of lattices. A frame change also
transports the neighboring edges and both endpoint tests.
Thus every summand of the transferred test has
\(r_j\le\dim V_j/g\) when \(4D_1<c\), by the previous
source rank argument.

\paragraph{Inputs to the dimension-weighted average.}
The dimension-weighted sum assumes rank bounds on its tensor constituents. We specify which sampled segment tests supply them on the enlarged box. For the enlarged weight box
below, the tensor edge test is required at depth \(D_0\)
and order \(\lfloor t_0(A,B)\rfloor\) with
\begin{equation}\label{eq:two-letter-tensor-rank}
 r_U\ge (1-\kappa_X(A,B))\dim U/g,\qquad
 \kappa_X(A,B)=\kappa_P+\sum_I\kappa_I(A,B).
\end{equation}
The sum is over segments, and \(\kappa_P\) accounts for the
digit blocks cut away by the margins. The two-letter construction supplies these tests at the base level.
Sections~\ref{sec:segment-window} and~\ref{sec:retained-rank} give \(\kappa_P\le C_s NN_1/Y\)
and, on the enlarged product box for fixed \(B\),
\[
 \mathbb E_A\kappa_I(A,B)
      \le C_A (s_I+1)N_1/Y,\qquad \sum_I(s_I+1)\le N.
\]
Here \(s_I+1\) is the number of sampled \(A\)-coordinates in
the segment, including its appended coordinate.

These coordinates have different indices in the sampling box, so their law is the
product marginal on that segment. In particular a uniform
bound for fixed \(B\) implies the same bound averaged over
both \(A\) and \(B\). This gives
\(\mathbb E\kappa_X\le(C_A+C_s)NN_1/Y\).

The function \(t_0\) is a constant independent of \(A,B\)
plus a sum of at most \(N\) terms of the form
\(\min(B_s,A_{s'}/p^e)\) with \(e\ge0\).
Let \(J=\lfloor J_0\rfloor\), where \(J_0\) is the cleaning
budget below. For \(0\le j_s\le J\), each term changes by
at most \(J_0\) under simultaneous subtraction of \(j\).
Thus the order
\[
 t_*(a',b')=\lceil t_0(a',b')+NJ_0\rceil
\]

for the lower weights is at least
\(\lfloor t_0(A,B)\rfloor\). Increasing the order only enlarges
the two endpoint lattices. Write \(F(\alpha,\beta)\) for
the least discarded fraction at depth \(D_1\) and order
\(t_*(\alpha,\beta)\) for the weight with these endpoints
and zero intermediate entries. When \(J\le\min_s(A_s,B_s)\),
the improved bound for the small \(j\)'s is
\begin{equation}\label{eq:transferred-small-deficit}
 r_j\le (1-F(A-j,B-j))\,\dim V_j/g.
\end{equation}
The comparison can now be stated entirely in terms of the
retained rank of the large tensor and the discarded fractions
of the transferred tests.

\begin{lemma}[Weighted rank deficit]\label{lem:transfer-rank-deficit}
Let \(J\le\min_s(A_s,B_s)\). Suppose the transferred ranks at
depth \(D_1\) satisfy \(\sum_j r_j\ge r_U\) and
\(r_j\le\dim V_j/g\) for every summand, with the sharper bound
\eqref{eq:transferred-small-deficit} for \(0\le j_s\le J\).
Assume the tensor lower bound \eqref{eq:two-letter-tensor-rank}. With the dimension weights of the indicated summands, the discarded fractions satisfy
\[
 \sum_{0\le j_s\le J\ (\forall s)}\frac{\dim V_j}{\dim U}F(A-j,B-j)\le\kappa_X(A,B).
\]

\end{lemma}

\begin{proof}
Using \eqref{eq:two-letter-tensor-rank} and
\(\sum_jr_j\ge r_U\) gives the weighted inequality
\begin{equation}\label{eq:edge-mean-weighted}
 \sum_{0\le j_s\le J} F(A-j,B-j)
    \prod_{s\in T}\frac{d(A_s-j_s,B_s-j_s)}{D(A_s)D(B_s)}
        \le \kappa_X(A,B).
\end{equation}
The upper rank bounds for the omitted summands are essential
to this inequality. The number \(g\) cancels before averaging.

\end{proof}

\begin{proposition}[From dimension weights to the product box]\label{prop:edge-box-deficit}
Let \(J=\lfloor J_0\rfloor\), \(e=J+c_{\rm cl}\) and
\(I_a^+,I_b^+\) be the respective upper enlargements by \(e\)
of the intervals \(I_a,I_b\) of lengths \(H_a,H_b\).
For \(0\le\tau\le c_{\rm cl}\), assume
\eqref{eq:edge-mean-weighted} on the enlarged box and the density
bound \eqref{eq:edge-marginal-density}, with \(\delta>0\), for every
coordinate pair in \((I_a+\tau)\times(I_b+\tau)\). Then
\[
 \mathbb E_{(I_a+\tau)^T\times(I_b+\tau)^T}F
 \le \delta^{-N}
    \left(\frac{(H_a+e)(H_b+e)}{H_aH_b}\right)^N
       \mathbb E_{(I_a^+)^T\times(I_b^+)^T}\kappa_X .
\]
The prefactor is less than two under \eqref{eq:edge-marginal-loss}. The assertion is uniform in the common shift \(\tau\).
\end{proposition}

\begin{proof}
Use integer intervals \(I_a,I_b\) of lengths \(H_a,H_b\)
for the original weight choice. Denote by \(I_a^+,I_b^+\)
their enlargements at the upper ends by \(e=J+c_{\rm cl}\),
where \(c_{\rm cl}\) bounds the permitted intermediate entries
and bounded shifts. For each \(0\le \tau\le c_{\rm cl}\)
and \((\alpha,\beta)\in(I_a+\tau)\times(I_b+\tau)\),
one has
\begin{equation}\label{eq:edge-marginal-density}
 \sum_{j=0}^J\frac{d(\alpha,\beta)}{D(\alpha+j)D(\beta+j)}
    \ge \delta,\qquad
 \delta=(1-(k-1)J_0/a)(1-\exp(-\Lambda_0/3)).
\end{equation}
Here \(a\) is the rounded lower endpoint of \(I_a\),
\(\Lambda_0=C_0N_1^2\log a\),
\(J_0=\Lambda_0 b_0/k\), and \(b_0\) is the reference
value for the second weight. Assume \(\alpha\ge a\),
\(\beta+J+1+(k-1)\le2b_0\), \(k\ge3\), and
\((k-1)J_0/a\le1\).

In fact
\(D(\alpha)/D(\alpha+j)\ge1-(k-1)J_0/a\)
by multiplying its reciprocal factors and applying Bernoulli's
inequality. Also
\[
 d(\alpha,\beta)/[D(\alpha)D(\beta)]\ge (k-1)/(\beta+k-1).
\]
Writing \(v_j=D(\beta)/D(\beta+j)\), the latter lower bound
times \(v_j\) dominates \(v_j-v_{j+1}\).
The final term \(v_{J+1}\) is at most
\(\exp(-(k-1)(J+1)/(2b_0))\le\exp(-\Lambda_0/3)\).
This proves the density for one top.

Write \(\mathcal B=(I_a+\tau)^T\times(I_b+\tau)^T\) and
\(\mathcal B^+=(I_a^+)^T\times(I_b^+)^T\). For every
\((\alpha,\beta)\in\mathcal B\) and \(0\le j_s\le J\), the source
\((A,B)=(\alpha+j,\beta+j)\) belongs to \(\mathcal B^+\), and each
such term occurs once for that \(j\) in the enlarged sum.
Keeping exactly those nonnegative terms in
\eqref{eq:edge-mean-weighted} gives
\[
 \sum_{(\alpha,\beta)\in\mathcal B} F(\alpha,\beta)
 \prod_{s\in T}\left(
 \sum_{j=0}^{J}\frac{d(\alpha_s,\beta_s)}
 {D(\alpha_s+j)D(\beta_s+j)}
 \right)
 \le
 \sum_{(A,B)\in\mathcal B^+}\kappa_X(A,B).
\]
This is a change of summation order over the product box and
uses no independence assumption on \(F\). Apply
\eqref{eq:edge-marginal-density} in each coordinate and divide
by \(\delta^N|\mathcal B|\) to obtain the stated prefactor. We use the
bounds below when
\begin{equation}\label{eq:edge-marginal-loss}
 N\left((k-1)J_0/a+\exp(-\Lambda_0/3)+e/H_a+e/H_b\right)
 \le 0.395 .
\end{equation}

\proofstep{Uniformity of the losses in the averaging.}
To check the constants, write \(\Lambda_*=C_0N_1^2\log a_{\rm real}\)
for the value before rounding, and
\(Z=\Lambda_* b_0/k\). The size hypotheses give
\(\Lambda_*\le\Lambda_0\le2\Lambda_*\), \(c_{\rm cl}\le Z\),
\(H_a,H_b\ge Y_bNZ/2\) and \(a\ge Y_b N\Lambda_* b_0\)
with \(Y_b\ge400\). These bound the first and last two
normalized terms by \(2/Y_b,6/Y_b,6/Y_b\).
Also
\[
 N\exp(-\Lambda_0/3)\le N(NY)^{-8/3}\le0.36
\]

for \(N\ge2,Y\ge1\).
Bernoulli bounds \(\delta^N\) below by \(1-s\), where
\(s=N((k-1)J_0/a+\exp(-\Lambda_0/3))\).
The enlargement factor is at most \(\exp(t)\), with
\(t=N(e/H_a+e/H_b)\). The numerical bounds give
\[
 e^t\le(1-t)^{-1},\qquad s+t\le0.395.
\]
Hence the prefactor
in the last mean is less than two.

\end{proof}

\subsection{The integral contraction for the intermediate weights}
\label{sec:edge-contraction}

The segment calculation uses weights with only the extreme entries varying. An integral contraction compares these tests to the prescribed weights with bounded intermediate entries. At one top put \(d=c_{\rm cl}\)
and \(r=k-2\). In this subsection write \(V_\nu\) for a Schur
space in \(k\)-row matrix polynomials with column degrees
\[
 \nu_0\ge\cdots\ge\nu_{k-1}\ge0.
\]
It consists of polynomials
transforming by \(\prod_i b_{ii}^{\nu_i}\) under right
multiplication of the matrix by an upper triangular matrix \(b\).
The flag minors have column sets consisting of initial columns.
The coordinate ring of \(\mathrm{GL}_k\) under left and right
translation identifies this space in characteristic zero with the
irreducible of highest weight \(\nu\).

The products of flag minors
of this shape span that irreducible, since their span is
row-stable and contains the highest vector. The integral standard
monomial theorem gives its basis over \(O\)
\cite[Theorem II.2.16]{abw-schur}. With a diagonal term order
on the flag minors the leading monomial of a standard tableau
records the entries of its rows. These monomials are distinct
because the tableau rows are weakly increasing. Thus the
span of the basis is saturated in the entire polynomial lattice.
For a weight \(\lambda\) whose last entry may be negative put
\[
 \nu=\lambda-\lambda_{k-1}{\bf 1}.
\]

The lattice in a frame with matrix \(g\) consists of
\[
 (\det g)^{-\lambda_{k-1}}f(g^{-1}X)
\]
with \(f\) in the
standard polynomial lattice. This is the determinant twist
used in \(H(L)\).

Take end weights \(A+d,-B-d\) and intermediate weight \(\mu\)
with \(|\mu_j|\le d\), and put
\begin{equation}\label{eq:edge-intermediate-weights}\begin{aligned}
 \lambda_-&=(A,0,\ldots,0,-B),&
 \lambda_+&=(A+2d,0,\ldots,0,-B-2d),\\
 S&=(d,\mu,-d),&
 S^*&=(d,-\mu_{\rm rev},-d).
\end{aligned}\end{equation}
We use a bilinear form on the \(r\)-row Schur lattices of
middle partitions
\(\alpha=B+d+\mu\) and \(\alpha'=d-\mu_{\rm rev}\).
Their entries satisfy
\(\alpha'_i=L-\alpha_{r+1-i}\) with \(L=B+2d\).
Put \(e=\alpha_r\) and \(\gamma=\alpha-e{\bf1}\).
Division by \(\det(Y)^e\) takes the first lattice to
\(V_\gamma(O)\), since each monomial of shape \(\alpha\)
has exactly \(e\) full-size minors.

For the second lattice,
with \(J_r\) the reversal matrix, use the map
\[
             f(Y)\longmapsto (\det Y)^{L-e} f(Y^{-\mathsf T}J_r).
\]
Jacobi's complementary-minor identity sends each initial-column
minor to an initial-column complementary minor divided by
\(\det Y\), up to sign. The determinant factor clears all
denominators and the output has shape \(\gamma\).
The same transformation in reverse, by its square on invertible
matrices, gives surjectivity on the integral lattice.

\begin{proposition}[Contraction across the intermediate weights]
\label{prop:integral-edge-contraction}
At each top take the weights of \eqref{eq:edge-intermediate-weights}
with \(|\mu_i|\le d\). The flat-weight Schur lattices admit
the multiplication and contraction maps through the integral
\(r\)-row Schur lattice of shape \(\gamma\) obtained from the
two complementary partitions \(\alpha,\alpha'\). For their tensor over \(N\) tops the exponent and the corresponding \(w\)-adic order are
\[
 d_T=Nd,\qquad \operatorname{ord}_w(p^{d_T})=2d_T\quad(w^2=p,\ \operatorname{ord}_w(w)=1).
\]
 After multiplication by \(p^{d_T}\), the pulled-back target
covector satisfies the target integral test and the normalized pairing
at the distinguished edge, with a loss of at most \(2d_T\) in the
\(w\)-adic conductor at the other incoming edges. The source test
holds simultaneously at all outgoing edges.
\end{proposition}

\begin{proof}
\proofstep{A uniform pairing for the intermediate representation.}
Here is a uniform integral pairing on \(V_\gamma\).
For the transpose adjoint, the rational contravariant dual has
the same highest weight as \(V_\gamma\). Choose the corresponding
isomorphism between the two rational modules to take the highest
polynomial to the highest dual functional evaluating to one
there. Both lattices have a primitive generator on that line
and are stable under divided powers. Write \(l=\gamma_1\)
and
\[
 h=\binom r2\lfloor\log_p(l+r-1)\rfloor
\]
for \(r\ge2\).
When \(r=1\) use \(h=0\).
In either lattice the submodule generated from the highest
vector by divided powers contains \(p^h\) times the full lattice.
Indeed take a basis \(u_j\) of the divided-power lattice in
that representation.

Apply the invariant Hom projector of the
preceding type calculation to \(u_j^\vee\otimes v\), where
\(v\) lies in the full lattice. The coordinate weights of
\(V_\gamma\) lie in \([0,l]\), as each symbol appears at most
once per column of a tableau. Each coordinate of a Hom weight is a difference of two such coordinates, so
\[
 0\le a,b\le l\quad\Longrightarrow\quad -l\le a-b\le l.
\]
Multiplying by \(p^h\) gives integral intertwiners,
which take the highest vector into the integral highest line
and its divided-power span into the divided-power span.
The trace identity for that projector reconstructs \(p^h v\)
as the sum of their values on \(u_j\), multiplied by
\(\dim V_\gamma\).

Let \(M\) and \(M^\sharp\) be the two full lattices and
\(L,L^\sharp\) the respective divided-power sublattices.
The contravariant isomorphism \(j\) sends \(L\) onto \(L^\sharp\).
The inclusions just proved in both lattices give
\[
 p^{2h}M^\sharp\subseteq p^h j(M)\subseteq M^\sharp .
\]
The right inclusion uses \(p^h M\subseteq L\), and the
left uses \(p^h M^\sharp\subseteq L^\sharp\).
Thus \(p^h j\) is an integral pairing map with cokernel
killed by \(p^{2h}\). Lifting \(p^{2h}\) times a dual basis
gives the integral partners with those diagonal values. This bound depends on the width
\(l\le2d\), even though the determinant power \(e\) can be
large. In particular for \(2d\le8a\) the depth is at most
\[
 \lceil (k-1)^2\log_p(8a+k-1)\rceil.
\]

Compose with determinant division and the complementary-minor
map to obtain a pairing \(J\) on shapes \(\alpha,\alpha'\),
with the same full set of partners. Simultaneous substitution
\(Y\mapsto gY\) in both inputs multiplies \(J\) by
\(\det(g)^L\). This follows from the transpose adjoint,
since the complementary-minor input then changes by
substitution by \(g^{-\mathsf T}\) and the factor
\(\det(g)^{L-e}\). The first input contributes the remaining
\(\det(g)^e\).

\proofstep{Contraction on the polynomial coefficients.}
For a \(k\)-row polynomial \(f(X)\), put
\(D(Y)=\operatorname{diag}(1,Y,1)\).
Apply the pairing \(J\) in the variable \(Y\), coefficientwise in \(X\), to define
\[
 \mathcal C(f,t)(X)=J_Y\bigl(f(XD(Y)),t(XD(Y))\bigr).
\]

Here the two \(k\)-row inputs have the normalized shapes
of \(\lambda=\lambda_-+S\) and \(S^*\).
Cauchy--Binet expands each flag minor in these substitutions
into middle flag minors of \(Y\) with coefficients in \(O[X]\)
for integral input. The \(Y\)-shapes are \(\alpha,\alpha'\)
because the first column of \(X\) is retained and the last is
unused. Thus contraction is integral. It commutes with left matrix substitution in \(X\).

For an invertible middle block \(g\) one has
\[
 D(g)D(Y)=D(gY),
\]
so the covariance of \(J\) gives
\[
 \mathcal C(f,t)(XD(g))=\det(g)^L\,\mathcal C(f,t)(X).
\]
This proves the middle diagonal degrees and invariance under
upper unipotent middle-column operations. A first-column
addition to a middle column is represented by a matrix \(u\)
whose only off-diagonal entry is in the first row. The matrix
\(D(Y)^{-1}uD(Y)\) is upper unipotent over \(K(Y)\), so
\[
 XuD(Y)=XD(Y)(D(Y)^{-1}uD(Y))
\]
has the same value in each input. Invariance there gives a
polynomial identity over \(K\). The last column is unused,
and the first-column degrees add. Together these operations
give precisely the right triangular character of the
normalized shape of \(\lambda_+\).

For any middle polynomial of shape \(d+\mu\), lift each
flag minor of size \(j\) to a flag minor including the first
row and first column, of size \(j+1\). If the product has \(a\) minors, the lifted product has first-row degree \(a\). Multiply by
\[
 X_{00}^{2d-a},\qquad 0\le a\le2d,
\]
to obtain first-row degree \(2d\).
Evaluation on \(D(Y)\) gives the original polynomial, and
the lift has shape \(S+d{\bf1}\) and last-row degree zero.
The same construction lifts the weight \(d-\mu_{\rm rev}\)
to the shape of \(S^*\). Since multiplication by
\((\det Y)^B\) identifies the shape \(d+\mu\) with \(\alpha\),
choose the lifted partners \(s_b,t_a\) so that
\[
 J((\det Y)^B s_b(D(Y)),t_a(D(Y)))
\]
is the chosen
scalar diagonal.

If \(F\) has the lower flat shape, its
substitution \(F(XD(Y))\) is \((\det Y)^B F(X)\),
by the first-column and cofactor formula.
It follows that \(\mathcal C(Fs_b,t_a)\) equals \(F\)
times a polynomial with the flat normalized shape
\((4d,2d,\ldots,2d,0)\), first-row degree \(4d\)
and last-row degree zero. This space is spanned by
\[
          Z=(X_{00}\det X_{[0,k-2],[0,k-2]})^{2d}.
\]
In its first-column and cofactor model, every variable factor
must have maximal first-row degree, and each cofactor must
omit the last row. Since \(Z(I_k)=1\), evaluation at the
identity matrix shows that
\[
 \mathcal C(Fs_b,t_a)= \delta_{ab}p^{2h}ZF.
\]

This constructs the one-component maps
\[
 i(F)=ZF,\ U_b(F)=s_b F,\ L_a(G)=\mathcal C(G,t_a).
\]

The number of pairs is \(\dim V_\gamma\), the dimension for
the middle weight. Tensor over all tops, reversing the two
end degrees at the odd tops, to obtain the maps on a cycle.

\proofstep{Verification on the edge lattices.}
We check the edge lattices for those maps. In a frame of an
edge the single changing row direction in the lattice has
scale \(1,w^{-1},w^{-2}\) at the source, intermediate
lattice and target. For first-row degree \(2d\) and last weight \(-d\) the determinant-twisted scaling factors are
\[
 w^d\quad\text{at the first change},\qquad w^{2d}\quad\text{at the second}.
\]
Thus the transported
\(s_b,t_a\) belong to the normalized intermediate lattice,
and \(w^{2d}t_a\) is integral at the target.
For \(Z\) the degree and the last weight are doubled.
It is primitive in the coefficient lattice of the intermediate
frame after normalization. To check saturation, reduce a product \(ZP\) in the coefficient
lattice modulo the uniformizer \(w\). Primitivity of \(Z\) and the
domain property give
\[
 \overline Z\,\overline P=0,\quad\overline Z\ne0
       \ \Longrightarrow\ \overline P=0 .
\]
The original coefficient lattices are saturated, so \(P/w\) then
belongs to the smaller lattice. This proves saturation of the image.

For any outgoing edge at the same source, a polynomial
integral at that source belongs to its normalized edge
lattice. Indeed in the respective adapted row coordinates
the transformed monomials have row weights between the
two end weights. In particular the normalization subtracts
the greatest one. This proves the outgoing tests for
\(i\) and \(U_b\), using integrality of multiplication
in every frame. For any incoming edge at the same target,
the target lattice for \(S^*\) likewise lies in its
normalized intermediate lattice. In the edge frame this uses
\[
 2j-(j-d)=j+d\ge0
\]
for every row weight \(j\in[-d,d]\).
For the tensor of \(t_a\)'s put \(d_T=Nd\).

At the target \(\mathcal C(\,\cdot\,,t_a)\) is integral after
multiplication by \(w^{2d_T}\), and it is integral at the
distinguished normalized edge with no multiplier.
At any other incoming normalized edge it loses at most
\(w^{2d_T}\), by the target inclusion just proved.
The normalized first weights of \(\lambda_+\) exceed those
of \(\lambda\) by \(d_T\). The pullback of a covector
multiplied by \(p^{d_T}\) therefore preserves the target
condition and the raw edge pairing, and decreases the
vanishing conductor at other incoming edges by at most
\(2d_T\). These are the properties needed in the next
subsection. They use the same integral section at all
neighbors of the fixed source or target.
\end{proof}


\subsection{Bounded changes of the full weight}

Apply the contraction to the two flat weights bracketing each prescribed top weight \((\alpha,\beta)\), with endpoints \((\alpha,\beta)\pm c_{\rm cl}\). Their deficits
are \(F_-\) and \(F_+\). The higher order \(t_*\) for the plus
weight can be used in both tests, since \(t_*\) is increasing in
each degree. The minus test is used when each shifted endpoint
is still in its original lower range. Write \(H_-\), \(H_+\)
and \(H_\lambda\) for the corresponding flat, enlarged flat,
and prescribed top weight spaces. Let \(d_-,d_+,d_\lambda\)
be their dimensions divided by \(g\).

The integral comparison needed here is stronger than a dimension
inequality. At each edge it uses a map \(i:H_-\to H_+\),
and maps \(U_b:H_-\to H_\lambda,\ L_a:H_\lambda\to H_+\),
such that
\[
 L_a U_b= \begin{cases}p^{q_b}i,&a=b,\\0,&a\ne b.\end{cases}
\]
The number \(b_0'\) of pairs of maps is the product, over the
tops, of the dimensions of the \(\mathrm{GL}_{k-2}\) modules
of the intermediate weights. The map \(i\) carries the
clean source lattice into the enlarged flat one and identifies
its normalized raw edge lattice with a saturated submodule
of the enlarged raw edge lattice.

Each \(U_b\) maps clean
sources to clean sources at the same order and conductor.
Pullback by each \(L_a\), with the correction
\(p^{A_1(H_+)-A_1(H_\lambda)}\), takes clean covectors on
\(H_+\) to clean covectors on \(H_\lambda\) with conductor
reduced by \(2Nc_{\rm cl}\). The denominator bound for these maps is
\(q_b\le N\lceil\kappa_1(p,a,k-1)\rceil\).
The contraction maps above have these properties.

\begin{proposition}[Edge tests with the restored bounded weights]\label{prop:bounded-weight-test-transfer}
Assume the maps \(i,U_b,L_a\) satisfying the indicated integral
properties, the two flat tests at depth \(D_1\) with deficits
\(F_-,F_+\), and \(4D_1<c\). On every nonboundary weight
where the two tests are available, denote the number of successful directions by \(r_\lambda\). In the prescribed size range,
\[
 r_\lambda\ge(1-F_--2F_+-\eta)d_\lambda,\qquad \eta=12Nkc_{\rm cl}/b_0.
\]

Under the following reserve and conductor bounds the depth is
at most \(D_p\) and the competitor conductor is \(c_\star\).
With the full loss assigned to the missing minus tests on the
lower boundary, the mean deficit over the original box is at most
\(6(C_A+C_s)NN_1/Y+2/Y\).
\end{proposition}

\begin{proof}
\proofstep{Recovering rank through the integral comparisons.}
Given these maps, the integral rank deduction goes as follows.
Use \(n_-,n_+\) successful directions in the respective flat
tests at depth \(D_1\), and make their two pairing matrices
scalar identities \(z_- I,z_+ I\). Transport the \(n_-\)
sources by \(i\), and pair them against the \(n_+\) clean
covectors in \(H_+\), giving a rectangular integral matrix
\(M\). Let \(r\) count its Smith factors dividing \(z_-z_+\),
and let \(s=n_--r\). In Smith coordinates the last \(s\) source rows therefore satisfy
\[
  M_{ij}\in w z_-z_+O\qquad(r<i\le n_-,\ 1\le j\le n_+).
\]
Here and below these are rows after the integral changes of bases.

Saturation gives an integral retraction of the raw edge
lattice in \(H_+\) onto that in \(H_-\), by Smith form.
Pull the \(n_-\) original covectors back along this retraction,
using the pairing scale \(p^{A_1(H_-)-A_1(H_+)}\).
They are integral on the raw edge lattice; they need no
additional vanishing at other edges. Take the column
combinations dual to the row operations of the preceding
paragraph and retain the \(s\) complementary columns.
Their evaluations on the \(s\) rows give \(z_- I_s\).
Now adjoin the \(n_+\) clean rows and covectors. The combined
matrix has blocks
\[
 \begin{pmatrix}z_- I_s & z_-z_+ E\\ B & z_+ I_{n_+}\end{pmatrix},
 \qquad E\in w\operatorname{Mat}(O),\quad B\in\operatorname{Mat}(O).
\]
The Schur complement of the second diagonal block is
\(z_- (I-EB)\).

For right-hand side \((a',b')\), solve explicitly:
\[
 x=(I-EB)^{-1}(z_-^{-1}a'-Eb'),\qquad
 y=z_+^{-1}(b'-Bx).
\]
Since \(E\) vanishes modulo \(w\), \(I-EB\) has an integral
inverse. Thus multiplication of the inverse of the combined
matrix by \(z_-z_+\) makes it integral. This gives \(s+n_+\)
directions at depth at most \(2D_1\), with clean source rows
and raw edge covectors. The outgoing-edge
rank bound used earlier needs extra vanishing only on the
sources, and therefore applies here when \(4D_1<c\).
Edge transitivity allows the same count at every outgoing edge.
It follows that
\[
 r\ge n_-+n_+-d_+.
\]


The low Smith factors of \(M\) supply \(r\) directions
with matrix \(z_-z_+ I_r\) using the clean covectors.
Apply \(U_b\) to these rows and the corrected pullbacks by
\(L_a\) to these covectors. The matrices are block diagonal
with blocks \(p^{q_b}z_-z_+ I_r\).
Scaling each block to the largest valuation, the test on
\(H_\lambda\) has at least
\[
 b_0'((1-F_-)d_- - F_+d_+)
\]
directions, at depth
\(2D_1+N\lceil\kappa_1(p,a,k-1)\rceil\).

\proofstep{Comparing the Weyl dimensions.}
The Weyl dimension formula compares these dimensions without
any loss that depends exponentially on the number of tops.
At a top write \(A=\alpha-c_{\rm cl}\), \(B=\beta-c_{\rm cl}\)
and \(\mu\) for the intermediate weight.
The internal \(\mu\)-factors form the \(\mathrm{GL}_{k-2}\)
dimension.

Each of the remaining \(2k-3\) factors changes by a
ratio between \(1\) and \(1+6c_{\rm cl}/b_0\) from the lower
flat weight to the prescribed weight, or from the lower flat
weight to the enlarged flat weight. This follows directly from
\(|\mu_i|\le c_{\rm cl}\) and \(A,B\ge b_0/3\). Hence
\[
 b_0' d_-\le d_\lambda\le {\cal E}b_0'd_-,
 \qquad d_+\le{\cal E}d_-,\qquad
 {\cal E}=(1+6c_{\rm cl}/b_0)^{(2k-3)N}.
\]
Put \(\eta=12Nk c_{\rm cl}/b_0\), with \(\eta\le1/Y\)
and \(\eta\le1/9\) in the size regime. Since
\({\cal E}\le\exp(\eta)\), one has \[
 1/{\cal E}\ge1-\eta,\qquad {\cal E}\le2.
\] Together with \(0\le F_-\le1\),
these estimates give at least
\[
 (1-F_--2F_+-\eta)d_\lambda
\]
directions.

At positions outside the tops all weights have absolute value
at most \(c_{\rm cl}\). Tensor with a pair of dual bases
in their normalized raw edge lattices. In an edge frame, the
three component lattices have a single coordinate at respective
scales \(1,w^{-1},w^{-2}\), so on the small-weight factors all
normalized orders differ by at most \(2n'c_{\rm cl}\) in
uniformizer units, where \(n'\) is the cycle length. In the
adapted basis let \(j\) be the total weight in that coordinate
and let \(a_S\) be the sum of the first weights of the small
factors. The two normalized exponents have opposite \(j\)-terms,
\[
 e_{\rm source}=j-a_S,\qquad e_{\rm dual}=-j-a_S.
\]

Their differences from the endpoint exponents \(0,-2j\)
are each between \(-2n'c_{\rm cl}\) and zero.
These same inclusions hold in the frames for all other
edges at either endpoint.
Thus the additional factors increase the endpoint order by
\(2n'c_{\rm cl}\), decrease the conductor by the same amount
and multiply the number of successful directions and the
total dimension by the same factor.
The conductor is first set to
\[
 c=c_*+2Nc_{\rm cl}+2n'c_{\rm cl}
\]
to allow these losses.

For the numerical comparison let \(a_r\) be the real lower
endpoint, \(a'=\lceil a_r\rceil\), \(L=\log_p a_r\) and
\(\ell=1+\log_p(8a'+k-1)\). The size hypotheses give
\(L\ge1\), \(\ell\le3L\), \(\log_p a'\le2L\), \(N_1^2\le k\),
\(2\le N\le k\) and \(Y\ge1\). Put \(P=NY^2k^3L\),
\(u_0=\lceil(k+1)\ell\rceil\), \(k_1=\lceil(k-1)^2\ell\rceil\)
and \(\kappa_0=2Nk\log_p a'\). There are \(n_s\le N\) segments.
Adding their depths and the two changes of lattices gives
\[
 D_1=n_s(C_A Y^2 N_1^2k^2\ell+2\kappa_0)+4\kappa_0+4Nu_0.
\]

In units of \(P\), the four terms on the right are bounded by
\(3C_A,8,16,16\), respectively, and \(Nk_1\le4P\).
Indeed \(u_0\le4kL\), \(k_1\le4k^2 L\), and
\(n_sN k\le Nk^2\). Thus the prescribed depth
\(D_p\ge C_H NY^2(N_1^5+k^3)L\), with \(C_H\ge8C_A+93\),
satisfies \(2D_1+Nk_1\le(6C_A+84)P\le D_p\).
As \(c_*\ge8(D_p+1)\), the strict inequalities used when
combining tests hold. The final conductor is exactly \(c_*\).

\proofstep{The depth and order after restoring the weights.}
For the order put \(\mathcal L=C_0N_1^2\log a_r\),
\(\mathcal L'=C_0N_1^2\log a'\), \(Q=p^{G'}\), and set
\[
 f(a,B)=a/p^2+b_0/p+Ba/(pk)+Bb_0/k.
\]
Here \(b_0\) is the lower reference parameter called \(b'\) in
the segment estimate. Let \(S\) be the sum over odd tops \(y\)
of \(\min(b_y,a_{\rm prev}/p^{d_{\rm prev}})\), with the
preceding even top and its distance taken cyclically. Each odd
top occurs once among all segment interiors. Under the plus
shift each summand increases by at most \(c_{\rm cl}\), giving
\(S'\le S+Nc_{\rm cl}\).

The pre-rounding order at these plus
degrees is \(t_0=S'+C_tN(f(a',\mathcal L')+kQ+Qc)\).
Adding \(N\mathcal L' b_0/k\), taking the ceiling and allowing
for the small factors gives
\begin{align*}
 t_*+2n'c_{\rm cl}\le{}&S+C NN_1 f(a_r,\mathcal L)+2n'c_{\rm cl}\\
 &+CN\big[Q(c_*+CNkL+4Nc_{\rm cl}+2n'c_{\rm cl})
                      +kQ+kL+c_{\rm cl}\big].
\end{align*}
Indeed \(a'\le2a_r\) and
\[
 \mathcal L'\le2\mathcal L,
\]
so the
four terms in \(f\) grow by factors at most \(2,1,4,2\).
The extra \(N\mathcal L'b_0/k\) uses the fourth term and the
ceiling costs at most \(1\le NkQ\). Taking \(C\ge4C_t+4\) proves
the bound, including the displayed conductor contributions.

\proofstep{Transport at the boundaries of the intervals.}
To transport a test, first identify the two degree-\(n'\)
unramified coefficient rings over \(\mathbb Z_p\). Their residue
field is \(\mathbb F_q\) with
\[
 q=p^{n'}.
\]
Each embeds in the
unramified ring of degree \(2n'\), with residue extension
\(\mathbb F_{q^2}/\mathbb F_q\). The specified isomorphism
on the smaller residue field extends to the larger ones, by
uniqueness of the quadratic finite field extension. Its Witt
lift extends the specified unramified-ring identification and
commutes with arithmetic Frobenius.

Adjoining the specified
root \(w^2=p\) to each larger ring gives the full valuation
ring by Eisenstein's criterion and total ramification. Extending
by
\[
 x+yw\mapsto\phi(x)+\phi(y)w'
\]
identifies both fraction
fields, the embedded smaller coefficient rings, the specified
uniformizers and the Frobenius on the smaller rings. Every
lattice membership and matrix pairing test therefore transports
in both directions, preserving its depth, conductor and integer
weight indices. Fixing one setting defines the same bound for
all settings at once. When there are no settings the uniform
assertion there is vacuous.

On the lower boundary assign \(F_-=1\) when a minus endpoint
falls outside its interval, and use the deficit bound
\(\kappa_{\rm text}=F_-+2F_++\eta\). For each nonboundary
point, shifting both families of degrees by \(-c_{\rm cl}\)
injects the point into the original box. Hence
\[
 \mathbb E F_-\le\mathbb E_{\rm box}F+
                     Nc_{\rm cl}(H_a^{-1}+H_b^{-1}),
 \qquad \mathbb E F_+=\mathbb E_{\rm plus\ box}F .
\]
The extra term bounds the union of the lower slices in the
\(2N\) intervals and is at most \(1/Y\); also \(\eta\le1/Y\).

Proposition~\ref{prop:edge-box-deficit} bounds both displayed
means of \(F\) by \(2(C_A+C_s)NN_1/Y\), using the same
enlargement and the fixed-\(B\) segment mean. Thus
\[
 \mathbb E\kappa_{\rm text}
 \le 6(C_A+C_s)NN_1/Y+2/Y .
\]
At \(Y=x^{13}\) this yields \(C N_vN_1/x^{13}\), where
\(N=N_v\), in Section~\ref{subsec:frobenius-selection}.
All allowances used for this estimate depend only on the
stated size data, rather than the edge or coefficient setting.

\end{proof}
